\section{公共基准评估}
\label{sec:experiemnt}

% 1. Baswed on what insights, we do the following...

% 2. Results (as expected..)

% \zixuan{some may move to introduction}

通过\cref{sec:analysis}中的受控实验与分析，我们发现：MAS 在子智能体能力边界上最有效；指令微调 LLM 比 RLM 更适合初始化编排器；MAS 在并行和对抗环境中能带来稳定收益。为了进一步验证 \ourframework 在合成环境之外的有效性与普适性，我们在一系列涵盖不同领域与任务特征的公共基准上对其进行评估。
% \zixuan{what practical guidence you get?}. 
% To futher demonstrate the effectivemnes of \ourframework, we conduct experiments on public benchmarks with various domains.

\textbf{数据集。}我们考虑五个常用基准，涵盖数学（\textbf{AIME24}、\textbf{AIME25}~\citep{aime2024}）、多跳问答
% \austin{Is GPQA multi-hop QA? I'd consider it more reasoning QA?} 
（\textbf{GPQA}~\cite{rein2023gpqagraduatelevelgoogleproofqa}、\textbf{HotpotQA}~\citep{DBLP:conf/emnlp/Yang0ZBCSM18}），以及基于多步搜索的问答（\textbf{BrowseComp+}~\cite{chen2025BrowseCompPlus}）。对于 AIME 和 GPQA，我们使用 \textbf{DeepScaleR}~\cite{deepscaler2025}作为训练数据。需要指出的是，DeepScaleR 主要关注数学推理，对 GPQA 而言基本可视为分布外（OOD）数据，因此 GPQA 结果反映 OOD 泛化性能。对于 HotpotQA，我们使用其给定的训练划分；对于 BrowseComp+，我们使用其中 80\% 的数据训练编排器。
% \shafiq{...}

\textbf{子智能体。}为保证公平比较，\ourframework 使用固定的候选子智能体集合；所有子智能体采用相同的 LLM 骨干，仅工具和提示工作流有所不同。
% \ryan{subagent pool is fixed, maybe worth saying explicitly}
我们纳入四种广泛使用的工作流来构成子智能体：\textbf{CoT}、CoT 自洽\textbf{（SC）}、\textbf{辩论}和\textbf{自我改进}。我们还加入一个\textbf{搜索}子智能体；它遵循 OpenDeepResearch\cite{langchain_open_deep_research}，以多轮方式自动决定何时搜索以及搜索什么。我们采用第~\ref{sec:analysis}节发现的最佳设置：以 \qwen 作为编排器，以 \ossonetwentylow 作为子智能体骨干。


\textbf{配置。}我们的分析表明，应根据子任务结构与底层 LLM 能力有选择地部署 MAS，而不应默认用它取代 SAS。所提出的 \textbf{\masness} 概念使 \ourframework 能显式编码这种依赖任务与模型的选择。鉴于 AIME 和 GPQA 中的数学、推理任务大多具有顺序性，我们采用\textbf{低 \masness}，至多允许从 \texttt{CoTAgent}、\texttt{SCAgent}、\texttt{DebateAgent} 和 \texttt{ReflexionAgent} 中选择一个子智能体。相比之下，HotpotQA 与 BrowseComp+ 包含更复杂的子任务结构（例如并行搜索），并且需要处在 SAS 能力边界上的能力（例如多轮信息检索）。因此，我们在这些基准上采用\textbf{高 \masness}，并额外加入 \texttt{DeepResearchAgent}。
% , with candidate sub-agents that additionally include a \texttt{DeepResearchAgent}. 
% \update{Specifically, for AIME24, AIME25, and GPQA, we use a low DoM setting with candidate sub-agents including CoT, SC, Debate, and Reflexion. For HotpotQA and BrowseComp+, we adopt a high DoM setting and additionally include a Search sub-agent.}
% \cref{tab:configs} summarizes the
% resulting configurations. 
% \ryan{was \masness for these datasets pre-specified or tuned?}
% Based on the understanding of these benchmarks (see Appendix~\ref{xxx} for more details), 
% we deploy the config below:
% \zixuan{use table}
% \begin{table}[t]
% \setlength{\fboxsep}{1pt}
% \centering
% \resizebox{\columnwidth}{!}{
% \begin{tabular}{ccc}
%\toprule
% \textbf{Benchmark} & \textbf{\masness} & \textbf{Candidate Sub-agents} \\
% \midrule
% AIME24 & Low & CoT, SC, Debate, Reflexion \\
% AIME25 & Low & CoT, SC, Debate, Reflexion \\
% GPQA & Low & CoT, SC, Debate, Reflexion \\
% HotpotQA & High & CoT, SC, Debate, Reflexion, Search \\
% BrowseComp+ & High & CoT, SC, Debate, Reflexion, Search \\
% \bottomrule
% \end{tabular}
% }
% \caption{\small \masness and candidate sub-agents for considered datasets.
% }
% \label{tab:configs}
% \end{table}

\textbf{总体结果。}\cref{tab:overall_public}给出了与相应 SAS 及其他可比 MAS 的对比结果，包括最先进（SoTA）的\textbf{推理时编排系统} \textbf{AFlow}\citep{zhang2024aflowautomatingagenticworkflow}、\textbf{MaAS}\citep{zhang2025multiagentarchitecturesearchagentic}、\textbf{MAS-Zero}\citep{Ke2025MASZero}，以及公开的 SoTA \textbf{训练时编排}系统 \textbf{MAS-GPT}\citep{ye2025masgpttrainingllmsbuild}和 \textbf{ToolOrchestra}\citep{su2025toolorchestraelevatingintelligenceefficient}（顺序式编排）。据我们所知，它们是仅有的发布了已训练编排器的训练时编排系统。为进行受控比较，推理时基线采用与 \ourframework 相同的编排器 LLM，训练时基线则采用各自官方发布的已训练编排器。所有系统都与 \ourframework 使用\textit{相同的候选子智能体和子智能体 LLM}，这可能不同于它们原本的训练环境（更多细节见第~\ref{ap:discuss_baselines}节和第~\ref{ap:addition_observation}节）。

% \shafiq{Give a brief description of those baselines}. 

% \zixuan{TODO: Note that AIME25 is generalized retutls}

% \begin{table}[h]
% % \setlength{\fboxsep}{1pt}
% \centering
% \resizebox{\columnwidth}{!}{
% \begin{tabular}{l cccc c}
% \toprule
% \multirow{2}{*}{\textbf{Method}}
% & \multicolumn{4}{c}{\textbf{IID Tasks}} 
% & \multicolumn{1}{c}{\textbf{OOD Task}} \\
% \cmidrule(lr){2-5}\cmidrule(lr){6-6}
% & \textbf{AIME24} & \textbf{AIME25} & \textbf{HotpotQA} & \textbf{BrowseComp+} & \textbf{GPQA} \\
% \midrule
% CoTAgent       & 50.00 & 45.00 & 33.56 & 1.12 & 60.54 \\
% SCAgent        & 57.50 & 51.67 & 35.50 & 0.75 & 62.88 \\
% DebateAgent    & 62.08 & 57.50 & 36.88 & 0.81 & 64.14 \\
% ReflexionAgent & 60.83 & 50.42 & 36.63 & 1.00 & 62.37 \\
% SearchAgent    & ---   & ---   & 46.44 & 8.56 & ---   \\
% \midrule
% \multicolumn{6}{l}{\textbf{Standalone Agents}} \\
% \midrule
% MaAS    & 32.50  &    &  &  &    \\
% MAS-Zero         & \multicolumn{5}{c}{
% No valid MAS generated with 7B orchestrator, aligned with \citet{Ke2025MASZero}
% } \\
% \midrule
% ToolOrchestra  &  22.92  &    &  &  &    \\
% \midrule
% \textbf{\ourframework} & \textbf{66.25} & \textbf{61.25} & \textbf{49.00} & \textbf{11.00} & \textbf{65.21} \\
% \bottomrule
% \end{tabular}
% }
% \caption{\small \avgat{8} of the considered benchmarks. ``---'' indicates non-applicable. 
% \zixuan{[Working 01/08] lacking baselines?, Training-based baselines include different sub-agents (HotpotQA, BrowseComp+) or datasets. We may only add baselines for AIME and \masness=low}
% % \zixuan{TODO: < 1\% can be removed}
% }
% \label{tab:overall_public}
% \end{table}


\begin{table}[h]
\centering
\setlength{\tabcolsep}{1.5pt}
\resizebox{\columnwidth}{!}{
\begin{tabular}{l cccc c}
\toprule
\multirow{2}{*}{\textbf{方法}}
& \multicolumn{4}{c}{\textbf{IID 任务}}
& \multicolumn{1}{c}{\textbf{OOD 任务}} \\
\cmidrule(lr){2-5}\cmidrule(lr){6-6}
& \textbf{AIME24} & \textbf{AIME25} & \textbf{HotpotQA} & \textbf{BrowseComp+} & \textbf{GPQA} \\
\midrule

\multicolumn{6}{c}{\textbf{独立智能体}} \\
\midrule
CoTAgent       & 50.00 & 45.00 & 33.56 & 1.12 & 60.54 \\
SCAgent        & 57.50 & 51.67 & 35.50 & 0.75 & 62.88 \\
DebateAgent    & 62.08 & 57.50 & 36.88 & 0.81 & 64.14 \\
ReflexionAgent & 60.83 & 50.42 & 36.63 & 1.00 & 62.37 \\
DeepResearchAgent    & ---   & ---   & 46.44 & 8.56 & ---   \\
\addlinespace[2pt]

\midrule
\multicolumn{6}{c}{\textbf{SoTA 推理时编排}} \\
\midrule
AFlow           & 62.50 & 53.33 & --- & --- & 65.43 \\
MaAS           & 32.50 & 40.83 & --- & --- & 40.78 \\
MAS-Zero       & \multicolumn{5}{c}{7B 编排器未生成有效 MAS} \\
\addlinespace[2pt]

\midrule
\multicolumn{6}{c}{\textbf{公开的 SoTA 训练时编排}} \\
\midrule
MAS-GPT   & 58.75 & 43.33 & ---   & ---   & 63.51 \\
ToolOrchestra   & 23.33 & 11.25 & 37.44 &  1.38  & 29.80 \\
\addlinespace[2pt]
\midrule
\multicolumn{6}{c}{\textbf{以 SoTA LLM 作为编排器}} \\
\midrule
GPT-5 & 55.00 &	47.72 &	25.87 &  0.50 & 59.01\\
Claude-Sonnet-4.5 & 45.56 &	35.00 & 38.00 &  0.50 & 21.72\\

\addlinespace[2pt]
\midrule
\multicolumn{6}{c}{\textbf{本文方法}} \\
\midrule
\textbf{\ourframework} & \textbf{66.25} & \textbf{61.25} & \textbf{49.00} & \textbf{11.00} & \textbf{65.21} \\
\bottomrule
\end{tabular}
}
\caption{\small 所考察基准上的 \avgat{8} 准确率。“---”表示不适用。
% \zixuan{4 baselines added. MAS-GPT has SC fallback, so it won't be too bad} 
% \shafiq{reduce the gaps between colomns}\zixuan{reduced}
% \zixuan{[Working 01/08] Lacking baselines? Training-based baselines include different sub-agents (HotpotQA, BrowseComp+) or datasets. We may only add baselines for AIME and \masness = low.}
}
\label{tab:overall_public}
\vspace{-\baselineskip}
\end{table}

% \begin{observation}
% \textbf{\ourframework consistently outperforms the underlying sub-agents {as well as SoTA inference-time and training-time orchestration systems} across all evaluated benchmarks, and it demonstrates strong OOD generalization on GPQA.} 
\textbf{\ourframework 在所有评估基准上始终优于强基线，并展现出稳健的 OOD 泛化能力。}
% \shafiq{update based on the newly added baselines}
% We see improvement and align with Sec. 4 observations
% \end{observation}
这些结果为 \ourframework 的有效性提供了具体证据，
% Across diverse benchmarks, the framework is able to leverage the
% capabilities of its sub-agents more effectively than any single sub-agent in isolation, 
表明性能提升来自编排，而不仅仅来自子智能体自身的能力。强劲的 OOD 泛化能力进一步说明，编排策略并未过拟合训练分布，而是能够迁移到未见问题。
% Moreover, the strong OOD generalization shows that the learned orchestration strategy does not overfit, but instead transfers to unseen problem distributions. 
综合来看，这些结果证实\cref{sec:analysis}中发现的原则可以推广到受控环境之外，并转化为真实基准上的稳定提升。

% also indicates that the findings in \cref{sec:analysis} is generalize beyond the \ourbenchmark. 

% MAS is helpful when parallel, and challenging sub-tasks are invovled. 

% \zixuan{TODO: explian what above ob means and how this show our effectiveness}


\textbf{进一步分析。}我们进一步考察 \ourframework 所提出 MAS 设计的结构。尽管\cref{ap:benchmark_generated_mas}给出了完整分析，我们仍在此强调几项关键发现，以说明学习到的编排如何适应不同任务。


\textbf{在低 \masness 下，\ourframework 学会了有效的单智能体委派。}
如\cref{fig:7b_120b_aime24_low}所示，\ourframework 学会了把整个任务委派给单个子智能体（20 步后委派率达到 100\%），并动态选择强子智能体，主要是 \texttt{ReflexionAgent} 与 \texttt{DebateAgent}；它们也是\cref{tab:overall_public}中表现最好的 SAS 基线。这表明，在低 DoM 下，\ourframework 能有效地把任务委派给最佳子智能体，从而取得更优性能。
% compared to baselines.


\textbf{在高 \masness 下，\ourframework 学会了利用并行性。}
如\cref{fig:7b_120b_browse_comp_plus_high}所示，\ourframework 学会调用 \texttt{DeepResearchAgent} 执行多路并行搜索，每个问题通常使用 3 到 4 个。由此，BrowseComp+ 上形成了一种典型 MAS 模式：先通过并行搜索收集相关信息，再由聚合智能体合并信息。正是这种全局编排策略带来了性能提升。
% {This reflects a typical MAS pattern produced by \ourframework on BrowseComp+: the system first retrieves potentially relevant information through multiple parallel searches and then introduces an aggregation agent to combine the retrieved results. This behavior is effective and highlights the global view adopted by \ourframework during MAS design: collecting comprehensive evidence before aggregation improves final answer correctness.}

% This indicate the typical pattern of generated MAS for BrowseComp+ first retrieves all potentially required information through multiple search agents and then introduces an organization or aggregation agent to combine the retrieved results. This behavior is effcive and clearly show the global view when \ourframework design the MAS: collecting comprehensive evidence before aggregation can improve final-answer correctness.

\textbf{\ourframework 达到了性能-成本帕累托前沿。}
如\cref{fig:parato_front}所示，在所考察的推理时与训练时编排系统中，\ourframework 位于帕累托前沿，以更低或相近的成本取得更高准确率。这种效率直接源于前述行为：\ourframework 会动态适应每个任务，生成符合底层子任务结构的 MAS 设计，并把执行委派给最有效的智能体配置。

% \textbf{\ourframework achieves perforamnce-cost parato-frnt. }As shwon in \cref{fig:parato_front}, \ourframework lies in the pareofron acorss the consdierd inference-time and training-time orchestation systems, achieving higher accuracy at lower or comparable cost. This is achieveable as \ourframework dynamically adapts to the given task by
% proposing MAS designs that align with the underlying sub-task structure and by
% delegating execution to the most effective agent configurations, as we have seen in the previoous observatiopn.

% \cref{tab:cost_breakdown}

% \zixuan{01/21 efficiency advantage? a subsec?}
% \zixuan{below pargraph is not neccessary}

% {Taken together, the observations under low and high \masness indicate that \ourframework dynamically adapts to the given task by
% proposing MAS designs that align with the underlying sub-task structure and by
% delegating execution to the most effective agent configurations.}



% \begin{discussion}
% \textbf{Training dynamic and Proposed MAS.}
% \zixuan{MISSING 6: statistic, dynamix and observation}

% \end{discussion}

% \begin{takeaway}
% \end{takeaway}
